\documentclass[11pt]{article}
\usepackage[margin=1in]{geometry}
\usepackage{amsmath,amssymb,amsthm,hyperref,enumitem}
\hypersetup{colorlinks=true,urlcolor=blue,linkcolor=blue,citecolor=blue}
\newtheorem{theorem}{Theorem}
\newtheorem{lemma}{Lemma}
\newcommand{\Q}{\mathbb Q}
\newcommand{\R}{\mathbb R}
\title{Uniform initial partial quotients\\of a square root of a zeta tail}
\author{Principia Math In Progress\\MathDB 375794}
\date{18 September 2026}
\begin{document}
\maketitle
\begin{abstract}
For $R_N=\sum_{j>N}j^{-3}$ and $g_N=\sqrt{R_N/2}$, we prove that
the first two partial quotients are $2N+1$ and $4N+2$ for every
$N\ge1$, and the third is $\lfloor(12N+6)/19\rfloor$ for every
$N\ge2$. This includes every Pell index in Feldman's Conjecture~4.1.
Every fixed partial quotient is eventually a quasi-polynomial.
The uniform claims follow from explicit rational telescoping bounds.
\end{abstract}
\noindent\textbf{Priority and campaign status.}
The mathematics below is complete. Campaign eligibility remains pending,
with zero new problem credit: Hu's 2026 dissertation explicitly claims
the general eventual quasi-polynomial mechanism, but its full text could
not be obtained to check whether it also contains these specific uniform
zeta-tail formulas. The general mechanism is not claimed as new here.
No Lean verification or independent review is asserted.

\section*{Task 1: the complete mathematical result}
Write $g_N=[0;a_1(N),a_2(N),\ldots]$ using regular continued fractions.
The irrationality of $\zeta(3)$, proved by Ap\'ery, implies that $g_N$
is irrational, so this expansion is infinite. Also $0<g_N<1$.
A function on the positive integers is eventually quasi-polynomial
if, for some period $M$, its restriction to each residue class modulo
$M$ agrees with a rational polynomial for all sufficiently large inputs.

\begin{theorem}\label{thm}
For every integer $N\ge1$,
\[
 a_1(N)=2N+1,\qquad a_2(N)=4N+2.
\]
For every integer $N\ge2$,
\[
 a_3(N)=\left\lfloor\frac{12N+6}{19}\right\rfloor.
\]
For each fixed $j\ge1$, $a_j(N)$ is eventually quasi-polynomial.
The first three eventual periods are $1,1,19$.
\end{theorem}
Here \emph{eventually} permits a threshold depending on $j$; no common
threshold for infinitely many coefficients is asserted. At $N=1$ the
third coefficient is $1$, whereas the displayed floor would be $0$.
Thus the word \emph{formal} in the source's pattern is essential.
All its Pell indices $N_k=(Q_{2k+1}-1)/2=3,20,119,\ldots$ are covered.
In particular the theorem proves both formulas singled out in
\cite[Conjecture~4.1, Remark~4.2, Open Problem~1]{F}.

\subsection*{1. Rational bounds for the tail}
Put $q=2N+1$ and define
\begin{align*}
 U_6(q)&=\frac2{q^2}-\frac2{q^4}+\frac{14}{3q^6},&
 L_8(q)&=U_6(q)-\frac{62}{3q^8},\\
 U_{10}(q)&=L_8(q)+\frac{762}{5q^{10}},&
 L_{12}(q)&=U_{10}(q)-\frac{5110}{3q^{12}}.
\end{align*}
For a rational function $f$ let
$D_f(q)=f(q)-f(q+2)-8/(q+1)^3$.
The exact coefficient certificates below show, for every $q>0$,
\[
 D_{U_6}>0,\quad D_{L_8}<0,\quad D_{U_{10}}>0,\quad D_{L_{12}}<0.
\]
Each of these functions tends to zero at infinity. Summing their
defects at $q,q+2,q+4,\ldots$ and using
$R_N=\sum_{r\ge0}8/(q+1+2r)^3$ therefore gives
\begin{equation}\label{bounds}
 \max(L_8(q),L_{12}(q))<R_N<\min(U_6(q),U_{10}(q)).
\end{equation}
All inequalities are strict.

\subsection*{2. The first two floors}
Let $u=1/g_N=\sqrt{2/R_N}$ and $\delta=u-q$.
The rational comparisons certified below are
\begin{align}
 L_8(q)&>\frac2{(q+1/(2q))^2} &&(q\ge3),\label{firstlo}\\
 U_6(q)&<\frac2{(q+1/(2q+1))^2} &&(q\ge5).\label{firsthi}
\end{align}
For $N=1$, monotonicity of $t^{-3}$ supplies the remaining upper bound:
\[
 R_1<\sum_{j=2}^{20}\frac1{j^3}+\int_{20}^{\infty}t^{-3}\,dt
 <\frac{1011}{5000}<\frac{49}{242}
 =\frac2{(3+1/7)^2}.
\]
The middle comparison is a finite rational inequality.
Consequently \eqref{bounds}--\eqref{firsthi}, including this small case,
give
\[
 \frac1{2q+1}<\delta<\frac1{2q}
 \quad\Longrightarrow\quad
 q<u<q+1,\qquad 2q<\frac1\delta<2q+1.
\]
These are exactly $a_1=q$ and $a_2=2q$.

\subsection*{3. The third floor}
The third complete quotient is
\[
 X_3=\frac1{1/\delta-2q}=\frac{\delta}{1-2q\delta}.
\]
This is strictly increasing in $\delta$ on the interval just proved.
Define
\[
 v_-(q)=q+\frac{6q}{12q^2+19},\qquad
 v_+(q)=q+\frac{6q^2+19}{q(12q^2+57)}.
\]
For $q\ge5$, the further rational comparisons below give
\[
 U_{10}(q)<\frac2{v_-(q)^2},
 \qquad L_{12}(q)>\frac2{v_+(q)^2}.
\]
Hence $v_-<u<v_+$, and substitution into the increasing expression
for $X_3$ yields the useful uniform estimate
\begin{equation}\label{x3}
 \frac{6q}{19}<X_3<\frac{6q}{19}+\frac1q\qquad(q\ge5).
\end{equation}
When $q\ge19$, the fractional part of $6q/19$ is a multiple of $1/19$,
and \eqref{x3} cannot cross the next integer. For the remaining odd
$q=5,7,9,11,13,15,17$, the exact positive gaps
$\lfloor6q/19\rfloor+1-6q/19-1/q$ are, respectively,
\[
 \frac{21}{95},\quad\frac{86}{133},\quad\frac8{171},\quad
 \frac{91}{209},\quad\frac{202}{247},\quad
 \frac{56}{285},\quad\frac{185}{323}.
\]
Thus $a_3=\lfloor6q/19\rfloor$ throughout $N\ge2$.
This floor is a quasi-polynomial of period $19$ in $N$.
Its least period is $19$: subtracting $12N/19$ leaves a bounded
periodic sequence, whose distinct fractional offsets require a period
divisible by $19$.

\subsection*{4. Every fixed coefficient is eventually quasi-polynomial}
We include the argument to make the application self-contained; the
general formal Gauss-map phenomenon is also stated in Hu's dissertation
abstract \cite{H}.
There is a unique formal series
$F(X)=\sum_{r\ge2}c_rX^{-r}$ satisfying
\begin{equation}\label{diff}
 F(X)-F(X+1)=(X+1)^{-3}.
\end{equation}
Its coefficients are rational and $c_2=1/2$: at each next order the
new coefficient $c_r$ appears multiplied by the nonzero integer $r$.
If $A_m(X)=\sum_{r=2}^m c_rX^{-r}$, then
\[
 A_m(N)-A_m(N+1)-(N+1)^{-3}=O(N^{-m-2}).
\]
Summing this identity to infinity proves
$R_N-A_m(N)=O(N^{-m-1})$. Hence $R_N$ has the full rational
asymptotic expansion $F$ and $g_N$ has the formal expansion
$G=\sqrt{F/2}\in\Q((X^{-1}))$, beginning with $1/(2X)$.
The square root with positive leading coefficient is used.

Neither $F$ nor $G$ is a rational function of $X$.
Indeed, if $F$ were rational, \eqref{diff} would be a rational-function
identity. It cannot be a polynomial, since it tends to zero and is
nonzero. Choose a pole $\alpha$ of maximal real part and a pole $\beta$
of minimal real part. The pole at $\alpha$ of $F(X)$ cannot be
cancelled by $F(X+1)$, so \eqref{diff} forces $\alpha=-1$.
Similarly the pole at $\beta-1$ of $F(X+1)$ cannot be cancelled by
$F(X)$, forcing $\beta=0$. These conclusions contradict the extremal
real parts. Since $F=2G^2$, rationality of $G$ is also impossible.

Proceed by induction through the ordinary floor-and-invert algorithm.
On each of finitely many residue classes, suppose all previous
coefficients are rational polynomials in $N$ for sufficiently large $N$.
The current complete quotient then has a full Laurent expansion
\[
 H(X)=\frac{A(X)G(X)+B(X)}{C(X)G(X)+D(X)},
 \qquad A,B,C,D\in\Q[X],\quad AD-BC=\pm1.
\]
It is nonrational, since otherwise the displayed equation would solve
for a rational $G$. Let $P(X)\in\Q[X]$ be its polynomial part.
The negative-power tail $H-P$ is nonzero; its first nonzero coefficient
determines its eventual sign, denoted $\sigma\in\{+,-\}$.
Let $d$ clear the denominators of all coefficients of $P$.
On each residue class modulo $d$, the fractional part of $P(N)$ is a
constant $\theta\in\{0,1/d,\ldots,(d-1)/d\}$, because
$P(N+d)-P(N)$ is an integer.
Refine the old residue classes by this period. Since the actual
complete quotient equals $P(N)$ plus a tail tending to zero,
its floor eventually equals
\[
 \begin{cases}
 P(N)-\theta,&\theta>0,\\
 P(N),&\theta=0,\ \sigma=+,\\
 P(N)-1,&\theta=0,\ \sigma=-.
 \end{cases}
\]
In each case this is a rational polynomial on the refined class.
Subtracting it and inverting gives a new full Laurent expansion.
The nonzero tail prevents formal termination, and the invertible
linear fractional transformation preserves nonrationality.
This proves the induction for every fixed depth, starting with
$H=1/G$. Only finitely many residue classes occur at any fixed depth,
so their finitely many thresholds have a common maximum.
This proves the final assertion of Theorem~\ref{thm}.
\hfill$\square$

\subsection*{5. Explicit algebra certificates}
For complete reproducibility define $P_d$ by the following coefficient
lists, ordered from degree $d$ to degree zero:
\begin{align*}
 P_6:\;&(124,744,1833,2372,1716,672,112),\\
 P_8:\;&(1143,9144,32676,68040,90280,78208,43200,13888,1984),\\
 P_{10}:\;&(76650,766500,3593323,10350584,20197784,27787328,\\
 &\hspace{12mm}27194720,18628928,8519104,2340864,292608),\\
 P_{12}:\;&(1414477,16973724,98558638,363216500,939352932,\\
 &\hspace{12mm}1785002880,2542240704,2722051008,2166240320,\\
 &\hspace{12mm}1245299712,489389568,117734400,13081600).
\end{align*}
Direct multiplication gives
\[
 D_{U_6}=\frac{8P_6}{3q^6(q+1)^3(q+2)^6},\qquad
 D_{L_8}=\frac{-8P_8}{3q^8(q+1)^3(q+2)^8},
\]
\[
 D_{U_{10}}=\frac{8P_{10}}{15q^{10}(q+1)^3(q+2)^{10}},\qquad
 D_{L_{12}}=\frac{-8P_{12}}{15q^{12}(q+1)^3(q+2)^{12}}.
\]
All listed coefficients are positive, proving the defect signs.
The four rational differences used in the proof have numerator $2B_i$
and the following positive denominators:
\[
\begin{array}{c|c|c}
i&\text{difference}&\text{denominator}\\ \hline
1&L_8-2/(q+1/(2q))^2&3q^8(2q^2+1)^2\\
2&2/(q+1/(2q+1))^2-U_6&3q^6(2q^2+q+1)^2\\
3&2/v_-^2-U_{10}&15q^{10}(12q^2+25)^2\\
4&L_{12}-2/v_+^2&15q^{12}(12q^4+63q^2+19)^2
\end{array}
\]
where
\begin{align*}
 B_1={}&19q^6-99q^4-117q^2-31,\\
 B_2={}&6q^5-16q^4-22q^3-32q^2-14q-7,\\
 B_3={}&10695q^6-93467q^4-588925q^2-714375,\\
 B_4={}&135q^{12}+54612q^{10}-718884q^8-14616460q^6\\
       &-53848988q^4-30170727q^2-4611775.
\end{align*}
Substitution $q=t+3$ in $B_1$ and $q=t+5$ in $B_2,B_3,B_4$
gives polynomials with strictly positive coefficients.
This is a finite coefficient test, not an asymptotic inference.
The accompanying exact verifier checks these identities and
coefficients, the seven floor gaps, and the $N=1$ bound; it also
independently encloses the first three coefficients for $1\le N\le18$
using rational tail bounds.

\section*{Task 2: effective thresholds at arbitrary depth}
Turn the residue-class induction into an algorithm that outputs an
explicit threshold $N_j$, constituent polynomials, and a period for
each fixed $a_j$. Quantify the precision needed in the tail enclosure.
The proof above establishes existence but does not bound $N_j$.

\section*{Task 3: period and degree growth}
Determine the least periods and polynomial degrees of later
coefficients. Can either grow faster than exponentially with the
depth? Identify which residue classes create unexpectedly large
partial quotients by cancellation.

\section*{Task 4: normalized tails of other odd zeta values}
For $s=2r+1$, study
$\sqrt{\bigl(\sum_{j>N}j^{-s}\bigr)/(2r)}$, whose leading term is
$1/(2rN^r)$. Derive its initial quasi-polynomials and uniform rational
enclosures, and compare the exceptional indices with the Pell
phenomenon at $s=3$.

\section*{Task 5: a formal proof of the uniform inequalities}
Formalize the telescoping bounds and four polynomial certificates,
then use the first two floor inequalities to discharge the source's
declaration \texttt{prop\_gcf}. This is a proposed follow-up:
the present package has not been checked by a proof assistant.

\begin{thebibliography}{9}
\bibitem{F} D.~V.~Feldman,
\emph{Anomalous Partial Quotients in the Continued Fraction of
$\sqrt{\zeta(3)-S_N}$}, arXiv:2607.04077v1 (5 July 2026),
\href{https://arxiv.org/abs/2607.04077}{arXiv record}.
Complete ten-page source inspected; Conjecture~4.1 is on p.~6.
\bibitem{H} Q.~Hu, \emph{Certain Hidden Stabilities of Continued
Fraction Arithmetic, with applications to classical constants and
functions of analysis}, Ph.D. dissertation, University of New Hampshire
(2026), \href{https://scholars.unh.edu/dissertation/2988/}{repository record}.
Only the complete public abstract and metadata were accessible;
the PDF returned HTTP 403. The abstract claims the general asymptotic
quasi-polynomial theorem, and is credited on that basis.
\bibitem{A} R.~Ap\'ery, \emph{Irrationalit\'e de $\zeta(2)$ et $\zeta(3)$},
Ast\'erisque \textbf{61} (1979), 11--13.
The classical irrationality theorem is used only to ensure that
each numerical continued fraction is infinite.
\end{thebibliography}
\end{document}
